% ---
% id: foundations/lem:enclosure-psd
% kind: lem
% title: "Positive semidefiniteness from an entrywise enclosure"
% summary: >-
%   A Hermitian matrix is at least μ times the identity when a Gaussian-rational
%   centre within entrywise radii r is at least ρ + μ times it, ρ the largest
%   row sum of r. Every hypothesis is a comparison of rationals, so matrices
%   with irrational entries are decided in exact arithmetic.
% prominence: standard
% status: implemented
% uses: []
% assumes: []
% refs: []
% provenance: newly-derived
% implements: certiqs_key.foundations.reference:enclosure_psd
% lean:
% ---

A matrix with irrational entries, such as one with entries
$S_{il}/\sqrt{n_in_l}$, is not a matrix of rationals, and an exact decision of
its positive semidefiniteness cannot read its entries directly. The lemma replaces the decision by one on a rational centre, at the
price of a margin. It is a sufficient condition only: a matrix whose least
eigenvalue is smaller than the uncertainty of its enclosure is in general not
certified, and a finer enclosure is the remedy.

\begin{lemma}\label{lem:enclosure-psd}
\begin{enumerate}[label=(\roman*),leftmargin=2.4em]
\item Let $X$ be a Hermitian complex $d\times d$ matrix, $X_c$ a Hermitian
  $d\times d$ matrix with entries in $\mathbb{Q}(i)$, and $r_{il}\ge0$
  rationals with $|X_{il}-(X_c)_{il}|\le r_{il}$ for all $i,l$. Let $\rho$ be
  a rational with
  \begin{equation}\label{eq:enclosure-radius}
    \rho\;\ge\;\max_{i}\sum_{l}r_{il}
    \qquad\text{or}\qquad
    \rho\ge0\ \text{ and }\ \rho^{2}\;\ge\;\sum_{i,l}r_{il}^{2},
  \end{equation}
  and $\mu\ge0$ a rational. If $X_c-(\rho+\mu)\,\idon{d}\succeq0$, then
  $X\succeq\mu\,\idon{d}$; in particular $X\succeq0$, and $X$ is positive
  definite when $\mu>0$. If $X_c-\rho\,\idon{d}$ is positive definite, then
  $X$ is positive definite. Every hypothesis on $X_c$, $r$, $\rho$ and $\mu$
  is a property of finitely many rationals and Gaussian rationals.
\item Let $n_1,\dots,n_d$ be positive rationals and
  $D\coloneqq\mathrm{diag}(\sqrt{n_1},\dots,\sqrt{n_d})$, and let $a_i,e_i$
  be rationals with $0\le e_i\le a_i$ and
  \begin{equation}\label{eq:sqrt-enclosure}
    (a_i-e_i)^{2}\,n_i\;\le\;1\;\le\;(a_i+e_i)^{2}\,n_i .
  \end{equation}
  Then $|a_i-n_i^{-1/2}|\le e_i$. Let $S$ and $Y$ be Hermitian $d\times d$
  matrices with entries in $\mathbb{Q}(i)$, put $X\coloneqq D^{-1}SD^{-1}-Y$,
  $A\coloneqq\mathrm{diag}(a_1,\dots,a_d)$, $X_c\coloneqq ASA-Y$ and
  \begin{equation}\label{eq:scaled-radii}
    r_{il}\;\coloneqq\;\sigma_{il}\,\big(a_ie_l+e_ia_l+e_ie_l\big),
    \qquad
    \sigma_{il}\coloneqq|\mathrm{Re}\,S_{il}|+|\mathrm{Im}\,S_{il}| .
  \end{equation}
  Then $X$, $X_c$ and $r$ satisfy the hypotheses of (i) on $X$, $X_c$ and
  $r$, and the radii do not depend on $Y$.
\item Let $X$, $X_c$ and $r$ be as in (i), let $F$ be a Hermitian positive
  semidefinite $d\times d$ matrix with entries in $\mathbb{Q}(i)$, and let
  $q\in[0,1]$ be rational. Then $X-(1-q)F$, $X_c-(1-q)F$ and $r$ satisfy the
  hypotheses of (i) on $X$, $X_c$ and $r$. If, for some $\rho$ and $\mu$ as
  in (i), $X_c-(1-q)F-(\rho+\mu)\idon{d}\succeq0$, then
  $X_c-(1-q')F-(\rho+\mu)\idon{d}\succeq0$ for every rational
  $q'\in[q,1]$.
\end{enumerate}
\end{lemma}

\begin{proof}
(i) Put $Z\coloneqq X-X_c$, a Hermitian matrix, so $|Z_{il}|=|Z_{li}|$. Then
$|Z_{il}|\le r'_{il}\coloneqq\min(r_{il},r_{li})$, and $r'$ is symmetric with
$\sum_lr'_{il}\le\sum_lr_{il}$ for every $i$. Let $x\in\mathbb{C}^{d}$ be a
unit vector. If $\rho\ge\max_i\sum_lr_{il}$, then, using
$|x_i||x_l|\le\tfrac12(|x_i|^{2}+|x_l|^{2})$ and the symmetry of $r'$,
\begin{equation*}
  |\langle x,Zx\rangle|
  \le\sum_{i,l}|x_i|\,r'_{il}\,|x_l|
  \le\tfrac12\sum_{i,l}r'_{il}\big(|x_i|^{2}+|x_l|^{2}\big)
  =\sum_i|x_i|^{2}\sum_lr'_{il}
  \le\max_i\sum_lr_{il}\le\rho .
\end{equation*}
If instead $\rho\ge0$ and $\rho^{2}\ge\sum_{i,l}r_{il}^{2}$, the
Cauchy--Schwarz inequality on the $d^{2}$ numbers $r_{il}$ and
$|x_i||x_l|$ gives
\begin{equation*}
  |\langle x,Zx\rangle|\le\sum_{i,l}r_{il}\,|x_i||x_l|
  \le\Big(\sum_{i,l}r_{il}^{2}\Big)^{1/2}\Big(\sum_{i,l}|x_i|^{2}|x_l|^{2}\Big)^{1/2}
  =\Big(\sum_{i,l}r_{il}^{2}\Big)^{1/2}\le\rho ,
\end{equation*}
since $\sum_{i,l}|x_i|^{2}|x_l|^{2}=\big(\sum_i|x_i|^{2}\big)^{2}=1$. In both
cases $\langle x,Xx\rangle=\langle x,X_cx\rangle+\langle x,Zx\rangle
\ge(\rho+\mu)-\rho=\mu$, because $X_c-(\rho+\mu)\idon{d}\succeq0$. As $X$ is
Hermitian and $x$ was an arbitrary unit vector, $X\succeq\mu\,\idon{d}$. If
instead $X_c-\rho\,\idon{d}$ is positive definite, the same bounds on
$|\langle x,Zx\rangle|$ give
$\langle x,Xx\rangle\ge\langle x,X_cx\rangle-\rho
=\langle x,(X_c-\rho\,\idon{d})x\rangle>0$ for every unit vector $x$, so $X$
is positive definite. The
hypotheses are a finite set of comparisons of rationals and the positive
semidefiniteness of the matrix $X_c-(\rho+\mu)\idon{d}$ with entries in
$\mathbb{Q}(i)$.

(ii) The map $t\mapsto t^{2}n_i$ is strictly increasing on $[0,\infty)$, and
$a_i-e_i\ge0$, $n_i^{-1/2}>0$, $(n_i^{-1/2})^{2}n_i=1$; so
\eqref{eq:sqrt-enclosure} gives $a_i-e_i\le n_i^{-1/2}\le a_i+e_i$. Write
$n_i^{-1/2}=a_i+\varepsilon_i$ with $|\varepsilon_i|\le e_i$. The entry
$(i,l)$ of $D^{-1}SD^{-1}$ is
$S_{il}(a_i+\varepsilon_i)(a_l+\varepsilon_l)$, and that of $ASA$ is
$S_{il}a_ia_l$, so
\begin{equation*}
  |X_{il}-(X_c)_{il}|
  =|S_{il}|\,\big|a_i\varepsilon_l+\varepsilon_ia_l+\varepsilon_i\varepsilon_l\big|
  \le\sigma_{il}\,\big(a_ie_l+e_ia_l+e_ie_l\big)=r_{il},
\end{equation*}
using $a_i,a_l\ge0$ and $|z|\le|\mathrm{Re}\,z|+|\mathrm{Im}\,z|$. The matrix
$X$ is Hermitian, $D^{-1}$ being real diagonal and $S$, $Y$ Hermitian; $X_c$
is Hermitian with entries in $\mathbb{Q}(i)$ for the same reasons; the $r_{il}$
are non-negative rationals, and $Y$ does not enter them.

(iii) $X-(1-q)F$ is Hermitian, and $X_c-(1-q)F$ is Hermitian with entries in
$\mathbb{Q}(i)$, $q$ being rational and $F$ having entries in
$\mathbb{Q}(i)$. Since
$\big(X-(1-q)F\big)_{il}-\big(X_c-(1-q)F\big)_{il}=X_{il}-(X_c)_{il}$, the
radii $r$ still bound the differences entry by entry. Finally,
$X_c-(1-q')F-(\rho+\mu)\idon{d}
=\big(X_c-(1-q)F-(\rho+\mu)\idon{d}\big)+(q'-q)F$ is a sum of two positive
semidefinite matrices, since $q'-q\ge0$ and $F\succeq0$.
\end{proof}
